\documentclass[../main.tex]{subfiles}
\begin{document}

\chapter{リオーダバッファ}
\lessoninfo{
  video={Lesson 8，動画 1--40},
  textbook={H\&P \cite{hennessy2017} §3.6，§3.8，§C.4；\cite{smith1988}},
  keywords={例外，正確な例外，ファントム例外，リオーダバッファ，コミットポインタ，発行ポインタ，コミット，分岐予測ミス回復，レジスタエイリアステーブル，統合リザベーションステーション，スーパースカラ，アウトオブオーダ実行},
}

第7章の Tomasulo のアルゴリズムは，各命令をそのオペランドが揃い次第実行
し，結果がブロードキャストされた時点でそれをレジスタファイルに書き込む．
リネーミングによって計算される値そのものは正しく保たれるが，実際の
プロセッサがどれも扱わなければならない2つの事象が，この仕組みを破綻
させる．例外と分岐予測ミスである．どちらの場合も，問題を起こした命令より
プログラム順で後ろにある命令がすでにレジスタを上書きしている可能性が
あり，その書き込みは取り消せない．この章では，命令をアウトオブオーダに
実行しながらレジスタの更新は厳密にプログラム順に行う，リオーダバッファを
導入する．そうして得られるプロセッサの動作を詳細に追い，最後に統合
リザベーションステーション，スーパースカラ実行，そしてアウトオブオーダ
プロセッサのどのステップが実際に順序どおりでないのかを正確に述べる．

\section{アウトオブオーダ実行における例外}
\label{sec:l08-exceptions-ooo}

例外とは，0 除算や不正なメモリアドレスへのアクセスのように，命令が実行
時に発生させる事象である．\source{Lesson 8，動画 1--2；H\&P §C.4，§3.6．}%
プロセッサは例外を起こした命令のアドレスを保存し，通常はオペレーティング
システムの一部である例外ハンドラに制御を移す．ハンドラが原因に対処し
終えると，保存されたアドレスに復帰することがあり，そのとき例外を起こした
命令が再び実行されて，プログラムはその後を続行する．ハンドラが正しく働く
ため，そして復帰後にプログラムが正しく続行するためには，レジスタとメモリ
が，命令を 1つずつ実行するプロセッサでの姿と同じように見えなければ
ならない．

\begin{definition}[正確な例外]
  \label{def:l08-precise}
  ハンドラが開始する時点で，例外を起こした命令よりプログラム順で先行する
  命令がすべて完了してレジスタとメモリを更新しており，かつ例外を起こした
  命令もそれより後続の命令もそれらを変更していないとき，その例外を
  \term[せいかくなれいがい]{正確な例外}[precise exception]という．
\end{definition}

Tomasulo のアルゴリズムは正確な例外を提供しない．次の例を考える．
\begin{lstlisting}[language={[x86masm]Assembler}, numbers=none]
I1:  DIV R1, R2, R3     ; R3 = 0: 0 除算
I2:  ADD R2, R4, R5     ; I1 が読む R2 を上書きする
I3:  ADD R6, R6, R7     ; R6 <- R6 + R7
\end{lstlisting}
I2 と I3 は I1 に依存しないので，発行から数サイクルのうちに実行され，R2 と
R6 に書き込む．除算は多数のサイクルを要し，ここで用いるモデルでは，除算器
は演算が完了した時点で 0 除算を報告する（\cref{fig:l08-exception}）．その
ときには後続の 2つの結果はすでにレジスタファイルに入っており，3つの不都合
が生じる．
\begin{itemize}
  \item ハンドラは，I1 で停止した逐次実行では生じえない R2 と R6 の値を
    見る．
  \item 復帰後，I1 が再び実行され，今度は I2 が書いた R2 の値を読む．
    プログラム順では決して見てはならない値である．
  \item I2 と I3 も 2 度目の実行を受ける．I2 は同じ値を計算し直すだけだが，
    I3 は R6 に R7 を 2 回加えてしまう．
\end{itemize}

\begin{figure}[htbp]
  \centering
  % Tomasulo's algorithm and an exception: two later instructions write their
% results into the register file long before the division reports the
% division by zero.
\begin{tikzpicture}[x=0.56cm,y=0.72cm, font=\footnotesize\sffamily, >={Stealth[length=1.8mm]},
  cell/.style={draw, fill=ln-box, minimum width=0.56cm, minimum height=0.44cm,
    inner sep=0pt, font=\scriptsize\sffamily},
  wr/.style={cell, draw=ln-caution, fill=ln-caution!15}]
  % cycle axis
  \foreach \c in {1,...,11} {
    \node[font=\scriptsize\sffamily, text=ln-muted] at (\c,0.68) {\c};
  }
  \node[anchor=east, font=\scriptsize\sffamily, text=ln-muted] at (0.4,0.68)
    {\iflangja{サイクル}{cycle}};
  % instruction labels
  \node[anchor=east] at (0.4,0)  {I1\enspace\texttt{DIV R1,R2,R3}};
  \node[anchor=east] at (0.4,-1) {I2\enspace\texttt{ADD R2,R4,R5}};
  \node[anchor=east] at (0.4,-2) {I3\enspace\texttt{ADD R6,R6,R7}};
  % I1: issue, long division, exception
  \node[cell] at (1,0) {I};
  \draw[fill=ln-box] (1.5,-0.22) rectangle (10.5,0.22);
  \node[font=\scriptsize\sffamily] at (6,0) {E\enspace(\iflangja{除算}{division})};
  \node[cell, draw=ln-caution, fill=ln-caution, text=white, font=\scriptsize\bfseries\sffamily]
    (exc) at (11,0) {!};
  % I2 and I3: short, independent
  \node[cell] at (2,-1) {I}; \node[cell] at (3,-1) {E}; \node[wr] (w2) at (4,-1) {W};
  \node[cell] at (3,-2) {I}; \node[cell] at (4,-2) {E}; \node[wr] (w3) at (5,-2) {W};
  % annotations
  \node[anchor=west, font=\scriptsize\sffamily, text=ln-caution] at (4.55,-1) {R2};
  \node[anchor=west, font=\scriptsize\sffamily, text=ln-caution] at (5.55,-2) {R6};
  \draw[ln-caution, thick, dashed] (11,-0.3) -- (11,-2.35);
  \node[anchor=north east, align=right, font=\scriptsize\sffamily, text=ln-caution]
    at (11.35,-2.35)
    {\iflangja{0 除算を報告：R2 と R6 はすでに上書き済み}%
              {division by zero reported: R2 and R6 already overwritten}};
\end{tikzpicture}

  \caption{Tomasulo のアルゴリズムのもとでの例外（I: 発行，E: 実行，
    W: 結果書き込み）．I2 と I3 は，除算 I1 が 0 除算を報告するよりはるか
    に早く R2 と R6 に書き込む．}
  \label{fig:l08-exception}
\end{figure}

\section{分岐予測ミス}
\label{sec:l08-mispredictions}

分岐予測ミスも同じ問題を引き起こす．\source{Lesson 8，動画 3；H\&P
§3.6．}次のプログラムでは，I2 が除算の結果を待っている間に，分岐予測器の
おかげでプロセッサは I3 から I6 までをフェッチし発行できる．
\begin{lstlisting}[language={[x86masm]Assembler}, numbers=none]
I1:        DIV R1, R2, R3     ; 多数のサイクルを要する
I2:        BEQ R1, R4, Skip   ; 不成立と予測される
I3:        ADD R5, R6, R7     ; I3-I6: 分岐が不成立の
I4:        LW  R8, 0(R5)      ;   ときだけ実行される
I5:        MUL R9, R5, R5
I6:        SUB R6, R9, R7
           ...
     Skip: ...
\end{lstlisting}
I2 は R1 を必要とするので除算が完了するまで実行できないが，I3 はただちに
オペランドが揃い，I4 もその 1 サイクル後に揃う．Tomasulo のアルゴリズムの
もとでは，R5 と R8 は分岐が解決されるよりはるかに早く書き込まれる．分岐が
成立すると判明した場合，誤った経路の命令はリザベーションステーションから
取り除き，\texttt{Skip} からフェッチをやり直せるが，R5 と R8 はプログラム
が決して計算しない値で上書きされており，正しい値は失われている．

第2の問題は，誤った経路の命令が例外を発生させるときに生じる．プログラムが
R5 をアドレスとして使うのは分岐が不成立のときだけであり，成立するときには
$\text{R6}+\text{R7}$ が有効なアドレスである必要はない．したがってロード
I4 は，I2 が実行されるよりはるかに早く例外を発生させうる．

\begin{definition}[ファントム例外]
  \label{def:l08-phantom}
  予測を外した経路上にあり，したがってプログラムの逐次実行では決して
  実行されない命令が発生させる例外を
  \term[ふぁんとむれいがい]{ファントム例外}[phantom exception]という．
\end{definition}

ファントム例外はハンドラに届いてはならない．プログラム順では，それを
発生させた命令は存在しないからである．

\section{正しいアウトオブオーダ実行}
\label{sec:l08-correct}

2つの問題の原因は同じである．命令が結果を計算し次第レジスタを変更して
しまい，先行するすべての命令が例外も予測ミスもなく完了したと判明する前に
書き込みが行われるのである．\source{Lesson 8，動画 4；H\&P §3.6；
\cite{smith1988}．}対策は，結果を計算することと，それを可視にすることとを
分離することである．
\begin{itemize}
  \item 命令はこれまでどおり，オペランドが揃い次第アウトオブオーダに
    \emph{実行}される．
  \item 結果もこれまでどおりアウトオブオーダに\emph{ブロードキャスト}され，
    待機中の命令は遅れなくそれを受け取る．
  \item しかし結果が\emph{レジスタに書き込まれる}のはプログラム順であり，
    しかもその命令がプログラムの実行に属すると確実になってからである．
\end{itemize}
書き込みが許されるまで結果を保持するには，発行された後の命令のプログラム
順を覚えておく構造が必要である．リザベーションステーションにはこれが
できない．命令は結果がブロードキャストされ次第ステーションを離れ，また
ステーションは特定の順序で割り当てられるわけでもないからである．

\begin{keypoint}
  実行はアウトオブオーダ，ブロードキャストもアウトオブオーダ，しかし
  レジスタへの書き込みはプログラム順に行う．
\end{keypoint}

\section{リオーダバッファ}
\label{sec:l08-rob}

プログラム順を覚えておく構造は，発行済みの命令のキューである．
\source{Lesson 8，動画 5--10；H\&P §3.6．}

\begin{definition}[リオーダバッファ]
  \label{def:l08-rob}
  \term[りおーだばっふぁ]{リオーダバッファ}[reorder buffer]%
  \index{ROB|see{リオーダバッファ}}（ROB）は，発行済みで，かつ結果を
  レジスタファイルに書き込む \cref{def:l08-commit} のコミットステップを
  まだ通過していない命令を，すべてプログラム順に保持する．各エントリは
  3つのフィールドを持つ．
  \begin{description}
    \item[Reg] 命令の宛先レジスタ．
    \item[Value] 計算された後の，その結果．
    \item[Done] Value が結果を保持していることを示すビット．
  \end{description}
\end{definition}

講義が描くのはこの 3つのフィールドだけである．実際にはエントリは，命令が
例外を発生させたか，あるいは予測を外した分岐であるかと，命令のアドレスも
記録しなければならない．コミットが例外を受け付けたりフェッチをやり直したり
するのに，これらが必要だからである（\cref{sec:l08-recovery}，
\cref{sec:l08-rob-exceptions}）．

ROB は 2つのポインタで管理される循環バッファである
（\cref{fig:l08-rob}）．\term[はっこうぽいんた]{発行ポインタ}[issue pointer]%
は次に発行される命令が受け取るエントリを示し，命令が発行されるたびに
1 エントリ進む．まだコミットしていない最も古い命令は
\term[こみっとぽいんた]{コミットポインタ}[commit pointer]の位置にある．
コミットポインタから発行ポインタの手前までのエントリが使用中であり，命令
をプログラム順に保持する．それ以外はすべて空きである．2つのポインタが
等しい状態は空のバッファと満杯のバッファの両方を表しうるので，使用中
エントリ数か満杯ビットで両者を区別する．通常の動作では両ポインタは進む
だけで，最後のエントリから先頭へ巻き戻る．発行ポインタが戻るのは回復の
ときだけである（\cref{sec:l08-recovery}）．すべてのエントリが使用中の
ときは，最も古い命令が ROB を離れるまで命令を発行できない．発行される
命令はすべてエントリを受け取り，レジスタに書き込まない分岐も例外では
ない．分岐のエントリはプログラム順における位置を保ち，後述するように予測
を外したかどうかを記録する．

\begin{figure}[htbp]
  \centering
  % The reorder buffer as a circular buffer: entries from the commit pointer
% up to the issue pointer are in use, in program order, and wrap around from
% the last entry to the first.
\begin{tikzpicture}[x=1cm,y=1cm, font=\footnotesize\sffamily, >={Stealth[length=1.8mm]},
  hd/.style={font=\footnotesize\bfseries\sffamily},
  used/.style={draw, fill=ln-accent!12, minimum height=0.45cm, inner sep=0pt},
  free/.style={draw, fill=white, minimum height=0.45cm, inner sep=0pt, text=ln-muted}]
  % column headings
  \node[hd] at (0.45,0.45) {Reg};
  \node[hd] at (1.55,0.45) {Value};
  \node[hd] at (2.65,0.45) {Done};
  % rows: entry / reg / value / done / used?
  \foreach \r/\reg/\val/\done/\st in {1/R2/{}/0/used, 2/R5/8/1/used,
      3/{}/{}/{}/free, 4/{}/{}/{}/free, 5/{}/{}/{}/free,
      6/R3/{}/0/used, 7/R1/17/1/used, 8/R3/25/1/used} {
    \pgfmathsetmacro{\y}{-(\r-1)*0.45}
    \node[\st, minimum width=0.9cm] at (0.45,\y) {\reg};
    \node[\st, minimum width=1.3cm] at (1.55,\y) {\val};
    \node[\st, minimum width=0.9cm] at (2.65,\y) {\done};
    \node[anchor=east, font=\scriptsize\sffamily, text=ln-muted] at (-0.05,\y) {ROB\r};
  }
  % pointers
  \draw[<-, thick, ln-accent] (3.15,-2.25) -- (4.0,-2.25)
    node[right, text=ln-accent] {\iflangja{コミットポインタ}{commit pointer}};
  \draw[<-, thick, ln-accent] (3.15,-0.9) -- (4.0,-0.9)
    node[right, text=ln-accent] {\iflangja{発行ポインタ}{issue pointer}};
  % wrap-around
  \draw[->, ln-muted] (-0.85,-3.15) -- (-1.15,-3.15) -- (-1.15,0.0) -- (-0.85,0.0);
  \node[rotate=90, font=\scriptsize\sffamily, text=ln-muted] at (-1.4,-1.57)
    {\iflangja{循環}{wrap around}};
  % program order of the entries in use
  \node[anchor=west, align=left, font=\scriptsize\sffamily, text=ln-muted] at (4.0,-3.15)
    {\iflangja{使用中のエントリ（プログラム順）：\\ROB6, ROB7, ROB8, ROB1, ROB2}%
              {entries in use, in program order:\\ROB6, ROB7, ROB8, ROB1, ROB2}};
\end{tikzpicture}

  \caption{8 エントリのリオーダバッファ．コミットポインタから発行ポインタ
    の手前までの網掛けのエントリが使用中である．ここではこの範囲が ROB8
    から ROB1 へ巻き戻っている．}
  \label{fig:l08-rob}
\end{figure}

\subsection{発行，ディスパッチ，結果書き込み}

ROB を備えると，Tomasulo のアルゴリズムの 3つのステップは次のように変わる．%
\index{とますろのあるごりずむ@Tomasulo のアルゴリズム}
\begin{description}
  \item[発行] 命令は，空いているリザベーションステーション（RS）と，発行
    ポインタの位置にある空き ROB エントリを必要とする．どちらかが欠ければ
    発行は待つ．ROB エントリは宛先レジスタを受け取り，Done ビットは
    クリアされ，発行ポインタが進む．続いて各オペランドを RAT で調べる．
    RAT のエントリは，いまや RS ではなく ROB エントリを示す．
    \begin{itemize}
      \item RAT エントリがレジスタファイルを指していれば，値をレジスタから
        複製する．
      \item Done ビットの立った ROB エントリを示していれば，値をその
        エントリから複製する．
      \item そうでなければ，RS はその ROB エントリをオペランドのタグとして
        記録し，そのタグを伴うブロードキャストを待つ．
    \end{itemize}
    最後に，宛先レジスタの RAT エントリを新しい ROB エントリに設定し，RS は
    そのエントリを自身の結果のタグとして記録する．
  \item[ディスパッチ] これまでと同様に，オペランドがすべて揃った命令が実行
    を開始する．RS はこの時点で解放される．
  \item[結果書き込み] 結果は，そのタグである ROB エントリとともに CDB に
    ブロードキャストされる．待機中の RS はこれまでどおり値を取り込み，ROB
    エントリは値を格納して Done ビットを立てる．レジスタファイルも RAT も
    変更されない．
\end{description}

2つ目の参照のケースは新しい．結果は結果書き込みからコミットまで ROB に
留まり，その間 RAT はその ROB エントリを示し続ける．この区間に発行される
命令は，値をそこで見つけなければならない．見つけられなければ
ブロードキャストを待つことになるが，そのブロードキャストはすでに
起きてしまっているからである．

RS をより早く解放する点も新しい．%
\index{りざべーしょんすてーしょん@リザベーションステーション}Tomasulo の
アルゴリズムでは，ステーション自体が結果をブロードキャストするときのタグ
なので，RS は結果書き込みまで確保され続けなければならない．ROB があれば
タグは ROB エントリなので，演算が始まってしまえば RS は不要になり，同じ
数の RS でより多くの待機中の命令を保持できる．

\subsection{コミット}

すべての命令を完了させる第4のステップがある．

\begin{definition}[コミット]
  \label{def:l08-commit}
  \term[こみっと]{コミット}[commit]\index{りたいあ@リタイア|see{コミット}}%
  のステップでは，プロセッサはコミットポインタの位置にあるエントリを
  調べる．その Done ビットが立っていれば，値を宛先レジスタに書き込み，その
  レジスタの RAT エントリがまだこの ROB エントリを示していれば，RAT
  エントリをレジスタファイルを指すように戻す．続いてコミットポインタが
  進み，エントリが解放される．
\end{definition}

コミットポインタの位置にある命令が完了していなければ，後続のエントリが
完了していても何もコミットされない．結果がレジスタファイルに入るのは
プログラム順にかぎられる．RAT エントリを戻すのは，それがまだコミットする
エントリを示している場合だけである．そうでなければ，後の命令が同じ
レジスタに書き込むということであり，RAT はその命令のエントリを示し続けな
ければならない（\cref{sec:l08-rat}）．コミットは\emph{リタイア}とも
呼ばれる．

\subsection{ハードウェア構成}

\cref{fig:l08-organization} は，第7章の構成に ROB を加えたものを示す．CDB
はいまや結果を RS と ROB に届け，RAT は ROB エントリを示し，レジスタ
ファイルに書き込むのはコミットだけである．命令は，フェッチ，デコード，
発行，実行（ディスパッチから演算の完了まで），結果書き込み，コミットという
6つのステップを経る．

\begin{figure}[htbp]
  \centering
  % Hardware organization of Tomasulo's algorithm with a reorder buffer: the RAT
% names ROB entries, the CDB delivers results to the RSs and the ROB, and only
% commit writes the register file.
\begin{tikzpicture}[x=1cm,y=1cm, font=\footnotesize\sffamily, >={Stealth[length=1.8mm]},
  row/.style={draw, fill=ln-box, minimum height=0.42cm, inner sep=0pt,
    font=\scriptsize\sffamily},
  lab/.style={font=\scriptsize\sffamily, text=ln-muted},
  hd/.style={font=\footnotesize\bfseries\sffamily},
  stepl/.style={font=\scriptsize\itshape\sffamily, text=ln-accent},
  unit/.style={draw, fill=white, trapezium, trapezium angle=65, shape border rotate=180,
    minimum width=1.3cm, minimum height=0.55cm, inner sep=1pt}]
  % ---------------------------------------------------------------- IQ
  \foreach \r in {1,...,4} {
    \node[row, minimum width=1.9cm] at (0.95,{-(\r-1)*0.42-0.21}) {};
  }
  \node[hd] at (0.95,0.25) {IQ};
  % ---------------------------------------------------------------- RAT
  \foreach \r/\v in {1/ROB2, 2/--, 3/ROB5, 4/ROB3} {
    \node[row, minimum width=0.95cm] at (3.45,{-(\r-1)*0.42-0.21}) {\v};
    \node[lab, anchor=east] at (2.95,{-(\r-1)*0.42-0.21}) {R\r};
  }
  \node[hd] at (3.45,0.25) {RAT};
  % ---------------------------------------------------------------- REGS
  \foreach \r/\v in {1/5, 2/-3, 3/12, 4/7} {
    \node[row, fill=white, minimum width=0.95cm] at (5.45,{-(\r-1)*0.42-0.21}) {\v};
    \node[lab, anchor=east] at (4.95,{-(\r-1)*0.42-0.21}) {R\r};
  }
  \node[hd] at (5.45,0.25) {\iflangja{レジスタ}{REGS}};
  % ---------------------------------------------------------------- ROB
  \node[lab] at (7.8,0.25) {Reg};
  \node[lab] at (8.8,0.25) {Value};
  \node[lab] at (9.8,0.25) {Done};
  \foreach \r/\reg/\val/\done/\f in {1/{}/{}/{}/white, 2/R1/9/1/ln-box, 3/R4/{}/0/ln-box,
      4/R3/{}/0/ln-box, 5/R3/21/1/ln-box, 6/{}/{}/{}/white} {
    \pgfmathsetmacro{\y}{-(\r-1)*0.42-0.21}
    \node[row, fill=\f, minimum width=0.8cm] at (7.8,\y) {\reg};
    \node[row, fill=\f, minimum width=1.2cm] at (8.8,\y) {\val};
    \node[row, fill=\f, minimum width=0.8cm] at (9.8,\y) {\done};
    \node[lab, anchor=west] at (10.25,\y) {ROB\r};
  }
  \node[hd, anchor=south] at (8.8,0.45) {ROB};
  \draw[->, thick, ln-accent] (7.38,-0.63) -- (6.65,-0.63) -- (6.65,-0.21) -- (5.97,-0.21);
  \node[stepl, anchor=north] at (6.65,-0.66) {\iflangja{コミット}{commit}};
  % ---------------------------------------------------------------- RSs
  \foreach \r in {1,2,3} {
    \node[row, minimum width=2.4cm] at (1.9,{-2.9-(\r-1)*0.42}) {};
    \node[lab, anchor=east] at (0.65,{-2.9-(\r-1)*0.42}) {AD\r};
  }
  \foreach \r in {1,2} {
    \node[row, minimum width=2.4cm] at (5.7,{-2.9-(\r-1)*0.42}) {};
    \node[lab, anchor=east] at (4.45,{-2.9-(\r-1)*0.42}) {MD\r};
  }
  % ---------------------------------------------------------------- units
  \node[unit] (add) at (1.9,-4.75) {$+,-$};
  \node[unit] (mul) at (5.7,-4.75) {$\times,\div$};
  \draw[->] (1.9,-3.95) -- (add.north);
  \draw[->] (5.7,-3.53) -- (mul.north);
  \node[stepl, anchor=west] at (1.95,-4.2) {\iflangja{ディスパッチ}{dispatch}};
  \node[stepl, anchor=west] at (5.75,-4.0) {\iflangja{ディスパッチ}{dispatch}};
  % ---------------------------------------------------------------- CDB
  \draw[line width=1.6pt, ln-accent] (1.9,-5.6) -- (8.8,-5.6);
  \draw[->] (add.south) -- (1.9,-5.6);
  \draw[->] (mul.south) -- (5.7,-5.6);
  \node[hd, text=ln-accent, anchor=north] at (4.0,-5.65) {CDB};
  \node[stepl, anchor=north west] at (5.8,-5.65) {\iflangja{結果書き込み}{write result}};
  % CDB to the RSs and to the ROB
  \draw[->, ln-accent] (3.55,-5.6) -- (3.55,-3.32) -- (3.12,-3.32);
  \draw[->, ln-accent] (7.35,-5.6) -- (7.35,-3.11) -- (6.92,-3.11);
  \draw[->, ln-accent] (8.8,-5.6) -- (8.8,-2.54);
  % ---------------------------------------------------------------- issue
  \draw[thick] (0.95,-1.7) -- (0.95,-2.31) -- (7.38,-2.31);
  \draw[->, thick] (7.0,-2.31) -- (7.38,-2.31);
  \draw[->, thick] (1.9,-2.31) -- (1.9,-2.68);
  \draw[->, thick] (5.7,-2.31) -- (5.7,-2.68);
  \draw[<->, dashed] (3.45,-1.7) -- (3.45,-2.31);
  \draw[->, dashed] (5.45,-1.7) -- (5.45,-2.31);
  \draw[->, dashed] (7.38,-1.89) -- (6.75,-1.89) -- (6.75,-2.31);
  \node[stepl, anchor=south east, inner xsep=1pt] at (7.36,-1.87)
    {\iflangja{Done なら値}{value if Done}};
  \node[stepl, anchor=south west] at (1.0,-2.31) {\iflangja{発行}{issue}};
\end{tikzpicture}

  \caption{リオーダバッファを備えた Tomasulo の構成．発行は RS と ROB
    エントリを確保し，RAT を通じてオペランドをリネームする．オペランドの値
    はレジスタファイルから，あるいは Done の立った ROB エントリから得られる．
    結果は CDB から RS と ROB へ送られ，コミットはコミットポインタの位置の
    値をレジスタファイルに複製する．}
  \label{fig:l08-organization}
\end{figure}

\begin{note}
  任意の時点で，これまでに発行された命令をプログラム順に見たときの，ある
  アーキテクチャレジスタの最新の値は，3つの場所のいずれかにある．RAT が
  レジスタファイルを指していればレジスタファイルの中，Done ビットの立った
  ROB エントリの中，あるいはまだ計算されておらず，Done ビットがクリア
  された ROB エントリに入るのを待っている状態である．レジスタファイル自体
  は，常に最後にコミットした命令の時点での値を保持する．
\end{note}

\section{分岐予測ミス回復}
\label{sec:l08-recovery}

ROB があれば，予測を外した経路の命令は，コミットする前に取り除かれる
かぎりレジスタファイルを変更できない．\source{Lesson 8，動画 11--12；
H\&P §3.6．}\cref{sec:l08-mispredictions} のプログラムでは，I2 が R1 を
待っている間に I3 から I6 が発行され，まだ実行中の I6 を除くすべてが結果を
自身の ROB エントリに書き込む．I2 がようやく実行され，分岐の予測が外れて
いたと分かると，予測ミスが I2 の ROB エントリに記録され，その時点では他に
何も起きない．分岐より前の命令は影響を受けず，実行とコミットを続ける．

分岐がコミットポインタに達したとき，それより前の命令はすべてコミットして
いるので，レジスタファイルは，逐次実行が分岐を実行し終えた時点で持つ値を
ちょうど保持している．プロセッサ内に残る命令はすべて分岐より後ろにあり，
誤った経路の上にある．そこでプロセッサは
\term[ぶんきよそくみすかいふく]{分岐予測ミス回復}[branch misprediction recovery]%
を行う．これは次の 4つの動作からなる（\cref{fig:l08-recovery}）．
\begin{enumerate}
  \item 発行ポインタをコミットポインタと等しく設定してすべての ROB
    エントリを破棄し，ROB を空にする．
  \item すべての RAT エントリを，レジスタファイルを指すように戻す．
  \item すべての RS を解放し，実行ユニット内の演算を放棄する．それらの結果
    は，破棄された ROB エントリにしか行き場がない．
  \item 分岐の正しい後続命令からフェッチをやり直す．
\end{enumerate}
予測が当たった分岐は単にコミットする．コミットポインタが進むだけで，他には
何も起きない．

\begin{figure}[htbp]
  \centering
  % Branch misprediction recovery when the mispredicted branch reaches the commit
% pointer: every later ROB entry is on the wrong path and is discarded, and
% every RAT entry is reset to the register file.
\begin{tikzpicture}[x=1cm,y=1cm, font=\footnotesize\sffamily, >={Stealth[length=1.8mm]},
  row/.style={draw, fill=ln-box, minimum height=0.45cm, inner sep=0pt},
  lab/.style={font=\scriptsize\sffamily, text=ln-muted},
  hd/.style={font=\scriptsize\bfseries\sffamily},
  note/.style={font=\scriptsize\sffamily, anchor=west}]
  % ROB
  \node[hd] at (1.25,0.42) {\iflangja{命令}{Instruction}};
  \node[hd] at (2.9,0.42) {Done};
  \foreach \r/\ins/\done/\f in {1/{}/{}/white, 2/{BEQ R1,R4,Skip}/1/ln-box,
      3/{ADD R5,R6,R7}/1/ln-caution!12, 4/{LW~~R8,0(R5)}/1/ln-caution!12,
      5/{MUL R9,R5,R5}/1/ln-caution!12, 6/{SUB R6,R9,R7}/0/ln-caution!12,
      7/{}/{}/white, 8/{}/{}/white} {
    \pgfmathsetmacro{\y}{-(\r-1)*0.45}
    \node[row, fill=\f, minimum width=2.5cm] at (1.25,\y) {\texttt{\ins}};
    \node[row, fill=\f, minimum width=0.8cm] at (2.9,\y) {\done};
    \node[lab, anchor=east] at (-0.05,\y) {ROB\r};
  }
  \node[note, text=ln-accent] at (3.35,-0.45) {\iflangja{予測ミス}{mispredicted}};
  \node[note, text=ln-caution] at (3.35,-1.35) {\iflangja{例外}{exception}};
  \node[note, text=ln-muted] at (3.35,-2.25) {\iflangja{実行中}{executing}};
  % pointers
  \draw[->, thick, ln-accent] (-1.45,-0.45) -- (-0.9,-0.45);
  \node[font=\scriptsize\sffamily, text=ln-accent, anchor=east] at (-1.45,-0.45)
    {\iflangja{コミット}{commit}};
  \draw[->, thick, ln-accent] (-1.45,-2.7) -- (-0.9,-2.7);
  \node[font=\scriptsize\sffamily, text=ln-accent, anchor=east] at (-1.45,-2.7)
    {\iflangja{発行}{issue}};
  % wrong path
  \draw[thick, ln-caution] (4.75,-0.72) -- (4.9,-0.72) -- (4.9,-2.43) -- (4.75,-2.43);
  \node[font=\scriptsize\sffamily, text=ln-caution, align=left, anchor=west] at (4.95,-1.57)
    {\iflangja{誤った経路：\\すべて破棄}{wrong path:\\all discarded}};
  % RAT before and after
  \node[hd] at (7.7,0.42) {\iflangja{回復前}{before}};
  \node[hd] at (8.7,0.42) {\iflangja{回復後}{after}};
  \node[hd, anchor=south] at (8.2,0.62) {RAT};
  \foreach \r/\reg/\before in {1/R1/--, 2/R5/ROB3, 3/R6/ROB6, 4/R8/ROB4, 5/R9/ROB5} {
    \pgfmathsetmacro{\y}{-(\r-1)*0.45}
    \node[lab, anchor=east] at (7.15,\y) {\reg};
    \node[row, minimum width=1.0cm, font=\scriptsize\sffamily] at (7.7,\y) {\before};
    \node[row, fill=white, minimum width=1.0cm, font=\scriptsize\sffamily] at (8.7,\y) {--};
  }
  \node[lab, align=left, anchor=north west] at (6.8,-2.2)
    {\iflangja{--：レジスタ\\\phantom{--：}ファイルを指す}{--: points to the\\\phantom{--: }register file}};
\end{tikzpicture}

  \caption{予測を外した分岐 I2 がコミットポインタに達したときの回復．I1 は
    すでにコミットしている．I3 から I6 の ROB エントリは，I4 に記録された
    例外もろとも破棄され，RAT のエントリはすべてレジスタファイルに戻る．}
  \label{fig:l08-recovery}
\end{figure}

\begin{note}
  コミット時に回復するのは最も単純な方式であり，講義で用いられるのもこれ
  である．実際のプロセッサは，予測を外した分岐が実行され次第回復を開始し，
  分岐より後ろの ROB エントリだけを破棄して直ちにフェッチをやり直す．
  こうすれば，分岐がコミットを待つ間にサイクルを失わずに済む（H\&P
  §3.6）．その場合 RAT は分岐の時点の状態に戻さなければならない．たとえば
  分岐の発行時に保存しておいた複製から戻す．
\end{note}

\section{リオーダバッファのもとでの例外}
\label{sec:l08-rob-exceptions}

例外は，それを発生させた命令のもう1つの結果として扱うことで処理する．
\source{Lesson 8，動画 13--15；H\&P §3.6．}命令が実行中に例外を発生させる
と，プロセッサはハンドラを呼び出す代わりに，その例外を命令の ROB エントリ
に記録する．そのエントリがコミットポインタに達したとき，例外が受け付け
られる．プロセッサは，予測ミス回復とまったく同様に ROB エントリを破棄し，
RAT を戻し，RS と実行ユニットを解放したうえで，例外を起こした命令の
アドレスを保存し，例外ハンドラのフェッチを開始する．

これにより \cref{sec:l08-exceptions-ooo}，\cref{sec:l08-mispredictions} の2つの
問題がともに解消される．
\begin{description}
  \item[正確な例外] 例外を起こした命令がコミットポインタに達した時点で，
    先行する命令はすべてコミットしており，後続の命令はレジスタファイルに
    書き込んでいない．ストアも同様にコミット時にはじめてメモリに書き込む
    （第9章）ので，メモリも後続の命令によって変更されていない．ハンドラは
    \cref{def:l08-precise} の状態をちょうど見ることになり，復帰後には後続の
    命令がその状態から実行される．
  \item[ファントム例外は起きない] 予測を外した経路上の例外は，予測を外した
    分岐より後ろのエントリに記録される．先に分岐がコミットし，その回復が
    そのエントリを例外もろとも破棄するので，例外が受け付けられることはない．
    \cref{fig:l08-recovery} では，I2 がコミットするときに I4 の例外が消える．
\end{description}

\subsection{実行の外部からの見え方}
\label{sec:l08-outside}

ROB は，命令がいつ実行されたのかという2つの捉え方を分離する．プロセッサの
内部では，命令は実行ユニットがその演算を行ったときに，順序を問わず実行
される．一方，プログラマ，オペレーティングシステム，例外ハンドラに見えるの
はレジスタファイル，メモリ，プログラムカウンタだけであり，これらにとって
命令が実行されるのはコミットしたときである．コミットはプログラム順に 1 命令
ずつ行われ，その様子は，各命令を完了させてから次を始めるプロセッサと
まったく変わらない．結果が ROB にあってもまだコミットしていない命令は，
外から見ればまだ実行されていない．その結果が破棄されれば，まったく実行
されなかったことになる．\cref{sec:l08-example} の例は，最後にその実例を示す．

\section{コミット時の RAT 更新}
\label{sec:l08-rat}

\cref{def:l08-commit} の RAT 更新に条件が付いているのは，1つのレジスタに
書き込む命令が同時に複数 ROB の中にありうるからである．\source{Lesson 8，
動画 16；H\&P §3.6．}\index{れじすたえいりあすてーぶる@レジスタエイリアステーブル}%
ROB2 と ROB5 の命令がともに R4 に書き込み，ROB5 がまだ ROB にある状態で
ROB2 がコミットするとする．R4 の RAT エントリは ROB5 を示している．そちらの
命令のほうが後で発行されたからである．
\begin{itemize}
  \item コミットは ROB2 の値を R4 に書き込む．これは ROB2 の後の
    アーキテクチャ上正しい値であり，ROB2 と ROB5 の間の命令が例外を発生
    させたときに例外ハンドラが見なければならない値でもある．
  \item それでも RAT エントリは ROB5 を示し続けなければならない．いま発行
    される命令はプログラム順で ROB5 より後ろにあり，レジスタファイルに入った
    ばかりの古い値ではなく，ROB5 の値を使わなければならない．
\end{itemize}
ROB5 がコミットするとき，その間にさらに後の書き込み命令が発行されていな
ければ，R4 の RAT エントリは依然として ROB5 を示しており，そのときにかぎって
エントリが戻される．逆に，コミット済みのエントリを示したまま RAT エントリを
放置することも，同じく誤りである．エントリはコミット時に解放され，
まもなく無関係な命令に与えられうるからである．

\section{リオーダバッファの例}
\label{sec:l08-example}

ここで，ROB を備えたプロセッサでプログラムを最後まで追う．\source{Lesson
8，動画 17--30；H\&P §3.6．}このプロセッサは，加算・減算用に 3つの RS
（AD1--AD3），乗算・除算用に 2つの RS（MD1，MD2），そして 8 エントリの ROB
を持つ．ADD と SUB は 1 サイクル，MUL は 4 サイクル，DIV は 6 サイクルを
要し，実行ユニットは十分にあって，準備のできた命令がユニットを待つことは
ない．各ステップは次の規約に従う．
\begin{itemize}
  \item 1 サイクルに 1 命令を発行し，必要な種類の空いている RS のうち番号の
    最も小さいものに入れる．
  \item サイクル $C$ のディスパッチで解放された RS と，サイクル $C$ の
    コミットで解放された ROB エントリは，サイクル $C+1$ から新しい命令に
    与えられる．
  \item 命令がディスパッチされるのは，発行の次のサイクル以降であり，かつ
    待っている各オペランドのブロードキャストの次のサイクル以降である．
  \item サイクル $D$ にディスパッチされ，実行時間が $k$ サイクルの演算は，
    サイクル $D$ から $D+k-1$ まで実行され，サイクル $D+k$ に結果を書き込む．
    CDB は 1本であり，この例では 2つの結果が同じサイクルに書き込まれること
    はない．
  \item 命令がコミットするのは早くとも結果書き込みの次のサイクルであり，
    1 サイクルにコミットする命令は高々 1つである．
\end{itemize}

\needspace{12\baselineskip}
\begin{example}
  \label{ex:l08-trace}
  プログラムは次のとおりである．
  \begin{lstlisting}[language={[x86masm]Assembler}, numbers=none]
I1:  DIV R1, R2, R3     ; R1 <- R2 / R3
I2:  ADD R4, R5, R6     ; R4 <- R5 + R6
I3:  MUL R2, R4, R4     ; R2 <- R4 * R4
I4:  SUB R5, R1, R6     ; R5 <- R1 - R6
I5:  ADD R4, R4, R3     ; R4 <- R4 + R3
I6:  ADD R6, R2, R4     ; R6 <- R2 + R4
\end{lstlisting}
  初期値は $\text{R1}=5$，$\text{R2}=48$，$\text{R3}=4$，$\text{R4}=2$，
  $\text{R5}=7$，$\text{R6}=3$ である．RAT のエントリはすべてレジスタ
  ファイルを指し，ROB の 2つのポインタはともに ROB1 にある．
\end{example}

\paragraph{サイクル 1--6：発行．}
1 サイクルに 1 命令が発行され，I$k$ は ROB$k$ に入る．I1 はレジスタ
ファイルから 48 と 4 を MD1 に複製し，R1 の RAT エントリを ROB1 に設定する．
サイクル~2 で I1 がディスパッチされて MD1 を解放し，I2 が値 7 と~3 とともに
AD1 に発行される．サイクル~3 で発行される I3 は，R4 の RAT エントリに ROB2
を見つける．ROB2 は完了していないので，I3 の両オペランドは ROB2 を待つ．
I3 は MD1 を再び使える．Tomasulo のアルゴリズムだけであれば，MD1 は
サイクル~8 まで I1 に占有されたままだったはずである．I4 は，サイクル~3 で
I2 が解放した AD1 に発行される．第 1 オペランドは ROB1 を待ち，第 2 オペランドは 3
である．サイクル~4 で I2 が結果を書き込み，ROB2 が 10 を受け取って I3 が
それを取り込む．サイクル~5 で I5 が AD2 に発行され，R4 の RAT エントリが
まだ ROB2 を示しているのを見つける．ROB2 は完了しているので，I5 は 10 を
ROB から複製する．古い値~2 をまだ保持しているレジスタファイルからでは
ない．続いて I5 は R4 を ROB5 にリネームする．サイクル~6 で I6 は AD3 に
発行される．同じサイクルに I5 のディスパッチで解放される AD2 は，
サイクル~7 からしか使えないからである．I6 は ROB3 と ROB5 を待つ．
\cref{tab:l08-rob-state,tab:l08-rat-state} はサイクル~6 の終了時の状態を
示す．AD1 が I4 を，AD3 が I6 を保持し，他の RS はすべて空きで，レジスタ
ファイルは変化していない．

\begin{table}[htbp]
  \centering\small
  \caption{\cref{ex:l08-trace} のサイクル 6 と 10 の終了時の ROB．
    サイクル~6 の終了時にはコミットポインタは ROB1 に，サイクル~10 の
    終了時には ROB3 にある．発行ポインタはどちらでも ROB7 にある．}
  \label{tab:l08-rob-state}
  \begin{tabular}{@{}llcrcrc@{}}
    \toprule
    & & & \multicolumn{2}{c}{サイクル 6} & \multicolumn{2}{c}{サイクル 10} \\
    \cmidrule(lr){4-5}\cmidrule(l){6-7}
    エントリ & 命令 & Reg & Value & Done & Value & Done \\
    \midrule
    ROB1 & \texttt{DIV R1, R2, R3} & R1 &     & 0 & \multicolumn{2}{c}{コミット済み} \\
    ROB2 & \texttt{ADD R4, R5, R6} & R4 & 10  & 1 & \multicolumn{2}{c}{コミット済み} \\
    ROB3 & \texttt{MUL R2, R4, R4} & R2 &     & 0 & 100 & 1 \\
    ROB4 & \texttt{SUB R5, R1, R6} & R5 &     & 0 & 9   & 1 \\
    ROB5 & \texttt{ADD R4, R4, R3} & R4 &     & 0 & 14  & 1 \\
    ROB6 & \texttt{ADD R6, R2, R4} & R6 &     & 0 &     & 0 \\
    \bottomrule
  \end{tabular}
\end{table}

\begin{table}[htbp]
  \centering\small
  \caption{\cref{ex:l08-trace} の RAT とレジスタファイル．ダッシュは，RAT
    エントリがレジスタファイルを指すことを表す．}
  \label{tab:l08-rat-state}
  \begin{tabular}{@{}lcrcrr@{}}
    \toprule
    & \multicolumn{2}{c}{サイクル 6} & \multicolumn{2}{c}{サイクル 10} & 最終 \\
    \cmidrule(lr){2-3}\cmidrule(lr){4-5}\cmidrule(l){6-6}
    レジスタ & RAT & 値 & RAT & 値 & 値 \\
    \midrule
    R1 & ROB1 & 5  & --   & 12 & 12  \\
    R2 & ROB3 & 48 & ROB3 & 48 & 100 \\
    R3 & --   & 4  & --   & 4  & 4   \\
    R4 & ROB5 & 2  & ROB5 & 10 & 14  \\
    R5 & ROB4 & 7  & ROB4 & 7  & 9   \\
    R6 & ROB6 & 3  & ROB6 & 3  & 114 \\
    \bottomrule
  \end{tabular}
\end{table}

\paragraph{サイクル 7--8：結果は出るがコミットはない．}
I5 はサイクル~7 に 14 を ROB5 に，I1 はサイクル~8 に 12 を ROB1 に書き込み，
I6 と I4 がこれらの値を取り込む．ROB2 はサイクル~4 から，ROB5 は
サイクル~7 から完了しているが，何もコミットされない．コミットポインタの
位置にある ROB1 が，サイクル~8 の終わりまで完了しないからである．

\paragraph{サイクル 9--10：最初のコミット．}
サイクル~9 で I1 がコミットする．R1 は 12 になり，R1 の RAT エントリはまだ
ROB1 を示しているので，レジスタファイルを指すように戻される．同じサイクル
に I3 が 100 を ROB3 に書き込み，I4 がディスパッチされる．サイクル~10 で
I2 がコミットして R4 は 10 になるが，R4 の RAT エントリは ROB5 を示している
ので，\cref{sec:l08-rat} の要求どおりそのままにされる．I4 が 9 を書き込み，
I6 がディスパッチされる．

\paragraph{サイクル 11--14：残りのコミット．}
I6 はサイクル~11 に 114 を書き込む．I3，I4，I5，I6 はサイクル 11 から 14 に
1 サイクル 1つずつコミットする．いずれも自身のレジスタへの最後の書き込み
命令なので，それぞれが自分の RAT エントリを戻す．最終的に ROB は空になり，
RAT のエントリはすべてレジスタファイルを指し，レジスタは $\text{R1}=12$，
$\text{R2}=100$，$\text{R3}=4$，$\text{R4}=14$，$\text{R5}=9$，
$\text{R6}=114$ という逐次実行の値を保持する．

サイクル~10 の終了時，レジスタファイルの R4 は 10 を保持する一方で，RAT は
R4 を読む命令を，14 を保持する ROB5 に導く．どちらも正しい．10 は最後に
コミットした命令 I2 の後の R4 の値であり，14 は最後に発行された書き込み
命令 I5 の後の値である．

この例は \cref{sec:l08-outside} の外部からの見え方も示している．I4 が
サイクル~9 の実行時に例外を発生させたとする．例外は ROB4 に記録され，I2 と
I3 はそのままサイクル 10 と~11 にコミットし，サイクル~12 に I4 がコミット
ポインタに達して例外が受け付けられる．ハンドラが見るのは $\text{R1}=12$，
$\text{R2}=100$，$\text{R4}=10$ であって，$\text{R4}=14$ ではない．I5 は
I1 が終わるより前のサイクル~7 に 14 を ROB5 に書き込んでいたにもかかわらず
である．外から見れば，実行はちょうど I4 の直前で停止したのである．

\section{リオーダバッファのもとでのタイミング}
\label{sec:l08-timing}

ROB を備えたプロセッサのサイクルごとの状態は，第7章の Tomasulo の
アルゴリズムと同様に，各命令が各ステップを通過するサイクルによって要約
でき，そこにコミットの列が 1つ加わる．\source{Lesson 8，動画 31--34；
H\&P §3.6．}記法は第7章のままとする．$I_n$，$D_n$，$W_n$ は命令 $n$ の
発行，ディスパッチ，結果書き込みのサイクル，$k_n$ はその実行時間，$P(n)$ は
そのオペランドを生成する命令の集合であり，$C_n$ をそのコミットのサイクルと
する．\cref{sec:l08-example} の規約のもとでは，$I_n$ は $I_{n-1}$ より後で，
必要な種類の RS と ROB エントリがともに空いている最初のサイクルである．
命令 $m$ のディスパッチで解放された RS はサイクル $D_m+1$ から，$m$ の
コミットで解放された ROB エントリはサイクル $C_m+1$ から空いている．実行
ユニットが十分にあり CDB の奪い合いもなければ，残りのサイクルは次式から
定まる．
\begin{equation}
  D_n = 1 + \max\Bigl(I_n,\; \max_{p \in P(n)} W_p\Bigr), \qquad
  W_n = D_n + k_n, \qquad
  C_n = \max\bigl(W_n + 1,\; C_{n-1} + 1\bigr),
  \label{eq:l08-timing}
\end{equation}
ここで第 1 式は，実行ユニットが空いている場合の \cref{eq:l07-dispatch} の
下限である．\tip{表はプログラム順に 1 行ずつ，各行は左から右へ，コミットを
最後にして埋めるとよい．}実行とは異なり，コミットは 2つの命令がデータを
共有しない場合でも直前の命令に依存する．RS をいつ解放するか，1 サイクルに
何命令コミットするかは設計上の仮定であり，これを変えれば表も変わる．第7章
のように RS が結果書き込みではじめて解放されるなら，RS はサイクル $W_m+1$
から空く．2つの命令が同じサイクルにコミットしうるなら，
$C_n = \max(W_n+1,\; C_{n-1},\; C_{n-2}+1)$ となる．

\begin{table}[htbp]
  \centering\small
  \caption{\cref{ex:l08-trace} のタイミング．実行の列は，演算が行われる
    サイクル $D_n$ から $D_n+k_n-1$ までを示す．}
  \label{tab:l08-timing}
  \begin{tabular}{@{}llcccc@{}}
    \toprule
    & 命令 & 発行 & 実行 & 結果書き込み & コミット \\
    \midrule
    I1 & \texttt{DIV R1, R2, R3} & 1 & 2--7 & 8  & 9  \\
    I2 & \texttt{ADD R4, R5, R6} & 2 & 3    & 4  & 10 \\
    I3 & \texttt{MUL R2, R4, R4} & 3 & 5--8 & 9  & 11 \\
    I4 & \texttt{SUB R5, R1, R6} & 4 & 9    & 10 & 12 \\
    I5 & \texttt{ADD R4, R4, R3} & 5 & 6    & 7  & 13 \\
    I6 & \texttt{ADD R6, R2, R4} & 6 & 10   & 11 & 14 \\
    \bottomrule
  \end{tabular}
\end{table}

\cref{tab:l08-timing} は例に対する結果を示す．実行の列と結果書き込みの列は
順序どおりでない．I5 は，先行する I1，I3，I4 より先に結果を書き込む．
コミットの列は増加している．I2 と I5 は結果書き込みの 6 サイクル後に
コミットする．ROB の先頭にある遅い除算に，そしてそれがサイクル~9 に
コミットした後は 1 サイクル 1 命令というコミットの上限に阻まれるからで
ある．それ以降の命令は，いずれも直前の命令のちょうど 1 サイクル後に
コミットする．この待ちはどの計算も遅らせない．I3 と I5 は I2 の結果を，
I6 は I3 と I5 の結果を，コミットを待つことなく得ている．ROB が遅らせる
のは，結果がアーキテクチャ状態の一部になる時点だけである．

\needspace{4\baselineskip}
とはいえ ROB は，小さすぎれば実行を制限しうる．エントリが 4つしかなければ，
I5 は I1 がコミットした次のサイクルであるサイクル~10 より前には発行できず，
オペランドがサイクル~5 から揃っているにもかかわらず，
\cref{tab:l08-timing} より数サイクル遅れて実行されることになる．したがって
ROB は，遅い命令がコミットを待つ間に発行されるすべての命令を保持できる
だけの大きさを持たなければならない．

\section{統合リザベーションステーション}
\label{sec:l08-unified}

\cref{fig:l08-organization} では，実行ユニットの種類ごとに専用の RS の組が
あり，各組の大きさは設計時に固定される．\source{Lesson 8，動画 35；H\&P
第3章．}プログラムによって演算の構成は異なるので，加算の多いプログラムで
は，乗算器の RS が空いているのに加算器の RS がすべて使用中になり，発行が
待たされることがある．命令はプログラム順に発行されるので，待たされている
加算より後ろにある乗算でさえ，空いている乗算器の RS を使えない．

\begin{definition}[統合リザベーションステーション]
  \label{def:l08-unified}
  すべての演算で共有される単一の RS の集まりを
  \term[とうごうりざべーしょんすてーしょん]{統合リザベーションステーション}[unified reservation station]%
  という．命令は，その集まりの中の空いている任意の RS に発行できる．
\end{definition}

このとき発行が待たされるのは，すべての RS が使用中のときだけであり，同じ数
の RS がより有効に使われる．代償はディスパッチにある．1つの実行ユニットの
ために 1つの組の RS から選ぶのではなく，ディスパッチ論理は，空いている実行
ユニットごとに，そのユニットが担当する演算の，準備のできた命令を集まり
全体から探し出し，その中から選ばなければならない．

\section{スーパースカラプロセッサ}
\label{sec:l08-superscalar}

第6章で見たように，プロセッサの IPC は，1 サイクルに扱える命令数によって
制限される（\cref{sec:l06-ipc}）．\source{Lesson 8，動画 36；H\&P §3.8．}

\begin{definition}[スーパースカラプロセッサ]
  \label{def:l08-superscalar}
  \term[すーぱーすから]{スーパースカラ}[superscalar]プロセッサは，すべての
  ステップを 1 サイクルに複数の命令に対して行う．すなわち，複数の命令を
  フェッチ，デコード，発行し，複数を実行ユニットにディスパッチし，複数の
  結果を書き込み，各サイクルに複数の命令をコミットする．
\end{definition}

そのためには，各ステップにそれぞれのハードウェアが必要になる．フェッチと
デコードは複数の命令を一度に扱わなければならない．発行は同じサイクルに
複数の命令をリネームし，そのそれぞれが自分より前の命令による RAT の更新を
見なければならない．ディスパッチには複数の実行ユニットが要る．複数の結果を
書き込むには複数の CDB が必要で，すべての RS のオペランドが自身のタグを
そのすべてと比較しなければならない．コミットは，コミットポインタの位置の
複数のエントリを調べ，複数のレジスタに書き込まなければならない．

これらのステップは，すべての命令が通る1本の鎖をなす．したがってプロセッサ
が持続できる 1 サイクル当たりの命令数は，最も制約の厳しいステップが扱える
数を超えない．1 サイクルに 4 命令をフェッチ，デコード，ディスパッチしながら
発行は 3 命令，コミットは 2 命令というプロセッサは，プログラムにどれだけ
ILP があっても，長い目で見れば 1 サイクルに高々 2 命令しか完了させられない．
1つのステップを広げることに意味があるのは，そのステップがボトルネックで
ある場合だけである．

\section{用語}
\label{sec:l08-terminology}

アウトオブオーダプロセッサのステップに対して本講義で用いる名称は，普遍的な
ものではない．\source{Lesson 8，動画 37；H\&P 第3章．}論文，教科書，
プロセッサの資料では，いくつもの別名が使われている
（\cref{tab:l08-terminology}）．\caution{同じ語が別のステップを指すことが
ある．文献によっては \emph{dispatch} が本講義でいう発行を，\emph{issue} が
本講義でいうディスパッチを指す．}その中にはコミットを \emph{complete} と
呼ぶものもある．これに対し本章では，演算が完了するのはその実行が終わった
ときである．

\begin{table}[htbp]
  \centering\small
  \caption{本講義と他の文献におけるステップの名称．}
  \label{tab:l08-terminology}
  \begin{tabular}{@{}ll@{}}
    \toprule
    本講義 & 他の文献での名称 \\
    \midrule
    発行 & issue，allocate，dispatch \\
    ディスパッチ & dispatch，execute，issue \\
    コミット & commit，complete，retire，graduate \\
    \bottomrule
  \end{tabular}
\end{table}

\section{何がアウトオブオーダなのか}
\label{sec:l08-order}

アウトオブオーダ実行\index{あうとおぶおーだじっこう@アウトオブオーダ実行}%
を行うプロセッサでも，すべてのステップが順序どおりでないわけではない
（\cref{fig:l08-order}）．\source{Lesson 8，動画 38--40；H\&P §3.6．}
\begin{description}
  \item[フェッチとデコード] 予測された経路に沿って，命令をプログラム順に
    取り込む．
  \item[発行] プログラム順である．リネーミングがそれを要求する．各命令は，
    自分より前のすべての命令が残した RAT を見なければならず，ROB エントリ
    もプログラム順に割り当てられるからである．
  \item[実行と結果書き込み] 順序どおりでない．命令はオペランドが揃い次第
    開始し，その演算が終わったときに終了する．後の命令が先の命令を追い越す
    のは，この2つのステップだけである．
  \item[コミット] 再びプログラム順である．実行が緩めた順序を ROB が回復
    する．
\end{description}

\begin{figure}[htbp]
  \centering
  % The steps of an out-of-order processor: only execution and write result
% are performed out of program order.
\begin{tikzpicture}[x=1cm,y=1cm, font=\scriptsize\sffamily, >={Stealth[length=1.6mm]},
  stp/.style={draw, fill=ln-box, minimum width=1.55cm, minimum height=0.62cm,
    inner sep=1pt, align=center},
  ooo/.style={stp, draw=ln-accent, fill=ln-accent!15}]
  \foreach \i/\txt/\sty in {0/{\iflangja{フェッチ}{Fetch}}/stp,
      1/{\iflangja{デコード}{Decode}}/stp, 2/{\iflangja{発行}{Issue}}/stp,
      3/{\iflangja{実行}{Execute}}/ooo, 4/{\iflangja{結果書き込み}{Write result}}/ooo,
      5/{\iflangja{コミット}{Commit}}/stp} {
    \node[\sty] (s\i) at ({\i*1.95},0) {\txt};
  }
  \foreach \i [evaluate=\i as \j using int(\i+1)] in {0,...,4} {
    \draw[->] (s\i.east) -- (s\j.west);
  }
  \draw[thick] ([yshift=4pt]s0.north west) -- ++(0,4pt) -| ([yshift=4pt]s2.north east);
  \draw[thick, ln-accent] ([yshift=4pt]s3.north west) -- ++(0,4pt) -| ([yshift=4pt]s4.north east);
  \draw[thick] ([yshift=4pt]s5.north west) -- ++(0,4pt) -| ([yshift=4pt]s5.north east);
  \node[anchor=base] at ([yshift=11pt]s1.north) {\iflangja{プログラム順}{in program order}};
  \node[anchor=base, text=ln-accent] at ([yshift=11pt]$(s3.north)!0.5!(s4.north)$)
    {\iflangja{アウトオブオーダ}{out of order}};
  \node[anchor=base] at ([yshift=11pt]s5.north) {\iflangja{プログラム順}{in program order}};
\end{tikzpicture}

  \caption{ROB を備えたアウトオブオーダプロセッサのステップ．プログラム順
    でないのは実行と結果書き込みだけである．}
  \label{fig:l08-order}
\end{figure}

\begin{keypoint}[章のまとめ]
  Tomasulo のアルゴリズムは，結果がブロードキャストされ次第それをレジスタ
  に書き込むので，例外は正確でなく，予測を外した命令がレジスタを壊し，
  誤った経路の命令がファントム例外を発生させうる．リオーダバッファは，
  発行済みの命令を，コミットポインタと発行ポインタの間にプログラム順に
  保持する．RAT は ROB エントリを示し，結果は RS と ROB にブロードキャスト
  され，RS はディスパッチ時に解放される．レジスタファイルに書き込み，
  コミットするエントリをまだ示している RAT エントリを戻すのは，プログラム
  順に ROB エントリを 1つずつ処理するコミットだけである．予測を
  外した分岐や例外を起こした命令への処置は，それがコミットするときに行う．
  ROB を空にし，RAT をレジスタファイルに戻し，RS と実行ユニットを解放して，
  正しいアドレスからフェッチをやり直す．タイミングの表にはコミットの列が
  加わり，これは \cref{eq:l08-timing} に従う．統合リザベーション
  ステーションは，1つの RS の集まりをすべての演算で共有する．スーパー
  スカラプロセッサはすべてのステップで 1 サイクルに複数の命令を扱い，最も
  狭いステップによって制限される．順序どおりでないのは，実行と結果書き込み
  だけである．
\end{keypoint}

\end{document}
